\chapter{Surfaces, Meshes, Point Clouds and Other Workbenches}

This chapter covers the workbenches you do not need every day but that save the day when it matters:
patching surfaces, repairing meshes, fitting point clouds, OpenSCAD, reverse engineering, inspection
and materials. They all live under \ui{More ▾}.

\section{Surface}

The Surface workbench has 7 commands, and they all work the same way: select the boundaries, create
the object, then tune the parameters in the dialog:

\begin{center}
\small
\begin{tabularx}{\linewidth}{@{}lX@{}}
\toprule
\textbf{Command} & \textbf{What it does} \\
\midrule
Filling & Surrounds a patch of surface with a set of boundary curves; you can set the degree (default 3), the number of points on the curves (15) and the number of iterations (2)\\
GeomFillSurface & Builds a transitional surface between 2, 3 or 4 boundaries; the fill type can be chosen\\
Sections & Builds a surface through several sections (like a loft, but it produces only a surface)\\
ExtendFace & Extends an existing face along the U/V directions (0.05× by default); can be symmetric\\
BlendCurve & A tangent transition curve between two edges\\
CurveOnMesh & Click control points on a mesh to project a spline curve onto it\\
Cut & Cuts two shapes against each other\\
\bottomrule
\end{tabularx}
\end{center}

A surface object is not a solid: thicken it with \ui{Part → Thickness} when you need to, or turn it
into a solid with \ui{Convert to solid} (after sewing).

\begin{practice}[Exercise: patching a hole in a shell]
Select the ring of edges around the hole → \ui{Filling} → adjust the degree until the surface is
smooth → \ui{OK} → select the surface and the shell → \ui{Part → Convert to solid}.
\end{practice}

\section{Mesh}

The Mesh workbench deals with triangle meshes: imported STL files, 3D scan results, and models
tessellated for analysis.

\subsection{Creating meshes}

\begin{itemize}
  \item \ui{Mesh from shape} (Standard): linear deflection 0.1, angular deflection 28.65°;
  \item \ui{Mefisto} mesher: subdivides by maximum edge length (default 1.0);
  \item \ui{Mesh From Geometry}: tessellates a shape;
  \item \ui{Regular solid}: cube, cylinder, cone, sphere, ellipsoid, torus;
  \item \ui{Remesh with Gmsh}: uses the built-in Gmsh (algorithm + maximum/minimum size + surface detection angle).
\end{itemize}

\subsection{Modifying and repairing}

Harmonizing/flipping normals, filling holes (by number of edges or by picking), adding a triangle
(three points), removing components, smoothing (Taubin/Laplace/median, 4 passes), decimating
(0.9 reduction), scaling, merging and splitting by components.

\subsection{Booleans and cutting}

Union/intersection/difference (at least two meshes, using the mesh kernel's own booleans); polygon
cut, trim, segment, cut by plane, cross-section (draw a polygon in the view first, then choose
whether to keep the inside or the outside); trim with a plane offset; sections (axis + number of
sections).

\subsection{Analyzing}

\ui{Mesh analysis} summarizes: number of points, number of faces, whether the mesh is a closed solid,
non-manifold edges, self-intersections and flipped faces, and offers one-tap repair buttons (Repair
all, remove duplicate points, fix degenerations, remove non-manifold, fix orientation, fix
self-intersections, remove wrinkles). There are also face evaluation, curvature, curvature info,
solid evaluation and bounding box.

\begin{practice}[Exercise: repairing an STL and converting it to a solid]
Import the STL → \ui{Mesh analysis} → look at the non-manifold edges/self-intersections → tap
\ui{Repair all} → \ui{Smooth} once → \ui{Part → Shape from mesh} (sewing tolerance 0.1) →
\ui{Convert to solid} → \ui{Refine shape} → export STEP.
\end{practice}

\section{Points}

\begin{itemize}
  \item \ui{Import points}: supports the usual point cloud text and binary formats;
  \item \ui{Export points}: \file{.asc} is recommended;
  \item \ui{Convert to points}: generates a point cloud from a shape or a mesh (point distance 0.1);
  \item \ui{Structured points}: arranges the point cloud into a regular grid (the column width can be set);
  \item \ui{Merge}: merges several point clouds into one;
  \item \ui{Poly cut}: draw a polygon in the view and keep the points outside or inside.
\end{itemize}

A point cloud is both the measured geometry for Inspection and the input for fitting surfaces in
reverse engineering.

\section{OpenSCAD}

The OpenSCAD workbench offers 15 commands, but \textbf{the app does not bundle an OpenSCAD
executable}, and you need to know this:

\begin{itemize}
  \item \textbf{Pure Python, works}: mirror/scale/resize mesh feature, refine shape, increase
        tolerance, edges to faces, expand placements, explode group, replace object, remove object
        and children, color shapes;
  \item \textbf{Needs an external binary, fails on mobile}: add OpenSCAD element, hull, Minkowski
        sum, mesh boolean. They report ``OpenSCAD executable unavailable''.
  \item The second toolbar reuses Part commands directly: check geometry, primitives, shape builder,
        cut/union/intersection, extrude, revolve.
\end{itemize}

\section{Reverse Engineering}

13 of the 15 commands work; the core idea is fitting geometry from points or meshes:

\begin{itemize}
  \item \ui{Approximate plane/sphere/cylinder}: least-squares fitting that outputs the
        corresponding surface;
  \item \ui{Approximate surface}: a B-spline surface; you can set the U/V degrees and the number of
        poles;
  \item \ui{Approximate curve}, \ui{Polynomial surface} (only reports the coefficients and the
        residual);
  \item \ui{Mesh boundary}: extracts a wire from the mesh boundary;
  \item \ui{Segmentation}: splits automatically by curvature tolerance (0.1) and minimum faces per
        segment (20);
  \item \ui{Segmentation from best-fit surfaces}: segments according to plane/cylinder/sphere fits;
  \item \ui{Manual segmentation}: draw polygons on the mesh to create segments;
  \item Missing: Poisson reconstruction and view triangulation.
\end{itemize}

\begin{practice}[Exercise: reverse engineering a scan]
Import the point cloud → \ui{Poly cut} to remove clutter → \ui{Approximate sphere} or
\ui{Approximate surface} → thicken the result or turn it into a solid → export.
\end{practice}

\section{Inspection}

Used for comparing surfaces:

\begin{itemize}
  \item \ui{Visual inspection}: select the measured geometry and the nominal geometry (multiple
        selection), set the search distance (default 0.05) and generate an inspection group; it
        paints a blue→green→red deviation map on the measured part;
  \item \ui{Inspection…}: tap the model to read the distance at that point; beyond the search radius
        it reports that there is no distance.
\end{itemize}

\section{Material}

\begin{itemize}
  \item \ui{Edit material}: opens the material library, where you edit the physical and appearance
        properties item by item; ``Save as new material'' puts it in the user library;
  \item \ui{Set material}: assigns a material to an object (and colors it according to its
        appearance);
  \item \ui{Inspect material}/\ui{Inspect appearance}: reads the current values;
  \item \ui{Color per face}: colors a single face;
  \item \ui{Texture mapping}: maps a texture onto an object.
\end{itemize}

\begin{tip}[Mass properties need a density]
\ui{Part → Mass properties} relies on the material's \textbf{density} to compute mass, center of
mass and moments of inertia. Set the material first, then compute the mass.
\end{tip}
